\documentclass[11pt]{report}
\usepackage[utf8]{inputenc}
\usepackage[spanish]{babel}
\usepackage{amsmath,amssymb,amsthm,mathtools}
\usepackage{xcolor}
\usepackage{listings}
\usepackage{booktabs}
\usepackage{hyperref}
\usepackage[T1]{fontenc}
\usepackage{bussproofs}
\EnableBpAbbreviations
\let\AXM\AXC \let\BIM\BIC \let\UIM\UIC \let\TIM\TIC \let\RLM\RL

% --- Notación ---
\newcommand{\Imp}{\rightarrow}
\newcommand{\Not}{\neg}
\newcommand{\Or}{\vee}
\providecommand{\Sat}{\models}
\providecommand{\Prov}{\vdash}
\newcommand{\Provnd}{\vdash_{\mathrm{ND}}}
\newcommand{\sem}[1]{[\![ #1 ]\!]}

% --- Entornos de teoremas (declaración del autor) ---
\newtheorem{teorema}{Teorema}[chapter]
\newtheorem{lema}[teorema]{Lema}
\newtheorem{corolario}[teorema]{Corolario}
\newtheorem{proposicion}[teorema]{Proposición}
\theoremstyle{definition}
\newtheorem{definicion}[teorema]{Definición}
\newtheorem{ejemplo}[teorema]{Ejemplo}
\theoremstyle{remark}
\newtheorem{nota}[teorema]{Nota}
\newenvironment{demostracion}{\begin{proof}[Demostración]}{\renewcommand{\qedsymbol}{$\blacksquare$}\end{proof}}
\newenvironment{solucion}{\begin{proof}[Solución]}{\renewcommand{\qedsymbol}{$\square$}\end{proof}}

% --- Configuración de listings para Lean ---
\lstdefinelanguage{Lean}{
  morekeywords={import, namespace, end, inductive, def, theorem, abbrev,
    by, intro, apply, exact, simp, constructor, cases, with, match, open,
    have, let, refine, rcases, induction, noncomputable, attribute,
    instance, obtain, by_cases, by_contra, fun},
  sensitive=true, morecomment=[l]{--}, morestring=[b]"}
\lstset{
  language=Lean, basicstyle=\ttfamily\small,
  keywordstyle=\color{blue!60!black}, commentstyle=\color{green!40!black},
  stringstyle=\color{red!50!black}, showstringspaces=false,
  breaklines=true, frame=single, columns=fullflexible, mathescape=true,
  literate=
    {⟶}{{$\to$}}1 {~}{{$\neg$}}1 {∅}{{$\emptyset$}}1 {⊢}{{$\vdash$}}1
    {⊧}{{$\models$}}1 {⊨}{{$\models$}}1 {⋎}{{$\lor$}}1 {∼}{{$\sim$}}1
    {¬}{{$\neg$}}1 {∀}{{$\forall$}}1 {→}{{$\to$}}1 {↔}{{$\leftrightarrow$}}1
    {Γ}{{$\Gamma$}}1 {Δ}{{$\Delta$}}1 {φ}{{$\varphi$}}1 {ψ}{{$\psi$}}1
    {χ}{{$\chi$}}1 {∈}{{$\in$}}1 {∉}{{$\notin$}}1 {⊆}{{$\subseteq$}}1
    {∪}{{$\cup$}}1 {≠}{{$\neq$}}1 {⊥}{{$\bot$}}1 {∧}{{$\land$}}1
    {∨}{{$\lor$}}1 {·}{{$\cdot$}}1 {⟨}{{$\langle$}}1 {⟩}{{$\rangle$}}1
    {▸}{{$\triangleright$}}1 {á}{{\'a}}1 {…}{{$\dots$}}1}

\title{Completitud de la deducción natural proposicional en Lean~4}
\author{}
\date{}

\begin{document}
\maketitle
\tableofcontents
\documentclass[11pt]{report}
\usepackage[utf8]{inputenc}
\usepackage[spanish]{babel}
\usepackage{amsmath,amssymb,amsthm,mathtools}
\usepackage{xcolor}
\usepackage{listings}
\usepackage{booktabs}
\usepackage{hyperref}
\usepackage[T1]{fontenc}
\usepackage{bussproofs}
\EnableBpAbbreviations
\let\AXM\AXC \let\BIM\BIC \let\UIM\UIC \let\TIM\TIC \let\RLM\RL

% --- Notación ---
\newcommand{\Imp}{\rightarrow}
\newcommand{\Not}{\neg}
\newcommand{\Or}{\vee}
\providecommand{\Sat}{\models}
\providecommand{\Prov}{\vdash}
\newcommand{\Provnd}{\vdash_{\mathrm{ND}}}
\newcommand{\sem}[1]{[\![ #1 ]\!]}

% --- Entornos de teoremas (declaración del autor) ---
\newtheorem{teorema}{Teorema}[chapter]
\newtheorem{lema}[teorema]{Lema}
\newtheorem{corolario}[teorema]{Corolario}
\newtheorem{proposicion}[teorema]{Proposición}
\theoremstyle{definition}
\newtheorem{definicion}[teorema]{Definición}
\newtheorem{ejemplo}[teorema]{Ejemplo}
\theoremstyle{remark}
\newtheorem{nota}[teorema]{Nota}
\newenvironment{demostracion}{\begin{proof}[Demostración]}{\renewcommand{\qedsymbol}{$\blacksquare$}\end{proof}}
\newenvironment{solucion}{\begin{proof}[Solución]}{\renewcommand{\qedsymbol}{$\square$}\end{proof}}

% --- Configuración de listings para Lean ---
\lstdefinelanguage{Lean}{
  morekeywords={import, namespace, end, inductive, def, theorem, abbrev,
    by, intro, apply, exact, simp, constructor, cases, with, match, open,
    have, let, refine, rcases, induction, noncomputable, attribute,
    instance, obtain, by_cases, by_contra, fun},
  sensitive=true, morecomment=[l]{--}, morestring=[b]"}
\lstset{
  language=Lean, basicstyle=\ttfamily\small,
  keywordstyle=\color{blue!60!black}, commentstyle=\color{green!40!black},
  stringstyle=\color{red!50!black}, showstringspaces=false,
  breaklines=true, frame=single, columns=fullflexible, mathescape=true,
  literate=
    {⟶}{{$\to$}}1 {~}{{$\neg$}}1 {∅}{{$\emptyset$}}1 {⊢}{{$\vdash$}}1
    {⊧}{{$\models$}}1 {⊨}{{$\models$}}1 {⋎}{{$\lor$}}1 {∼}{{$\sim$}}1
    {¬}{{$\neg$}}1 {∀}{{$\forall$}}1 {→}{{$\to$}}1 {↔}{{$\leftrightarrow$}}1
    {Γ}{{$\Gamma$}}1 {Δ}{{$\Delta$}}1 {φ}{{$\varphi$}}1 {ψ}{{$\psi$}}1
    {χ}{{$\chi$}}1 {∈}{{$\in$}}1 {∉}{{$\notin$}}1 {⊆}{{$\subseteq$}}1
    {∪}{{$\cup$}}1 {≠}{{$\neq$}}1 {⊥}{{$\bot$}}1 {∧}{{$\land$}}1
    {∨}{{$\lor$}}1 {·}{{$\cdot$}}1 {⟨}{{$\langle$}}1 {⟩}{{$\rangle$}}1
    {▸}{{$\triangleright$}}1 {á}{{\'a}}1 {…}{{$\dots$}}1}

\title{Completitud de la deducción natural proposicional en Lean~4}
\author{}
\date{}

\begin{document}
\maketitle
\tableofcontents
\documentclass[11pt]{report}
\usepackage[utf8]{inputenc}
\usepackage[spanish]{babel}
\usepackage{amsmath,amssymb,amsthm,mathtools}
\usepackage{xcolor}
\usepackage{listings}
\usepackage{booktabs}
\usepackage{hyperref}
\usepackage[T1]{fontenc}
\usepackage{bussproofs}
\EnableBpAbbreviations
\let\AXM\AXC \let\BIM\BIC \let\UIM\UIC \let\TIM\TIC \let\RLM\RL

% --- Notación ---
\newcommand{\Imp}{\rightarrow}
\newcommand{\Not}{\neg}
\newcommand{\Or}{\vee}
\providecommand{\Sat}{\models}
\providecommand{\Prov}{\vdash}
\newcommand{\Provnd}{\vdash_{\mathrm{ND}}}
\newcommand{\sem}[1]{[\![ #1 ]\!]}

% --- Entornos de teoremas (declaración del autor) ---
\newtheorem{teorema}{Teorema}[chapter]
\newtheorem{lema}[teorema]{Lema}
\newtheorem{corolario}[teorema]{Corolario}
\newtheorem{proposicion}[teorema]{Proposición}
\theoremstyle{definition}
\newtheorem{definicion}[teorema]{Definición}
\newtheorem{ejemplo}[teorema]{Ejemplo}
\theoremstyle{remark}
\newtheorem{nota}[teorema]{Nota}
\newenvironment{demostracion}{\begin{proof}[Demostración]}{\renewcommand{\qedsymbol}{$\blacksquare$}\end{proof}}
\newenvironment{solucion}{\begin{proof}[Solución]}{\renewcommand{\qedsymbol}{$\square$}\end{proof}}

% --- Configuración de listings para Lean ---
\lstdefinelanguage{Lean}{
  morekeywords={import, namespace, end, inductive, def, theorem, abbrev,
    by, intro, apply, exact, simp, constructor, cases, with, match, open,
    have, let, refine, rcases, induction, noncomputable, attribute,
    instance, obtain, by_cases, by_contra, fun},
  sensitive=true, morecomment=[l]{--}, morestring=[b]"}
\lstset{
  language=Lean, basicstyle=\ttfamily\small,
  keywordstyle=\color{blue!60!black}, commentstyle=\color{green!40!black},
  stringstyle=\color{red!50!black}, showstringspaces=false,
  breaklines=true, frame=single, columns=fullflexible, mathescape=true,
  literate=
    {⟶}{{$\to$}}1 {~}{{$\neg$}}1 {∅}{{$\emptyset$}}1 {⊢}{{$\vdash$}}1
    {⊧}{{$\models$}}1 {⊨}{{$\models$}}1 {⋎}{{$\lor$}}1 {∼}{{$\sim$}}1
    {¬}{{$\neg$}}1 {∀}{{$\forall$}}1 {→}{{$\to$}}1 {↔}{{$\leftrightarrow$}}1
    {Γ}{{$\Gamma$}}1 {Δ}{{$\Delta$}}1 {φ}{{$\varphi$}}1 {ψ}{{$\psi$}}1
    {χ}{{$\chi$}}1 {∈}{{$\in$}}1 {∉}{{$\notin$}}1 {⊆}{{$\subseteq$}}1
    {∪}{{$\cup$}}1 {≠}{{$\neq$}}1 {⊥}{{$\bot$}}1 {∧}{{$\land$}}1
    {∨}{{$\lor$}}1 {·}{{$\cdot$}}1 {⟨}{{$\langle$}}1 {⟩}{{$\rangle$}}1
    {▸}{{$\triangleright$}}1 {á}{{\'a}}1 {…}{{$\dots$}}1}

\title{Completitud de la deducción natural proposicional en Lean~4}
\author{}
\date{}

\begin{document}
\maketitle
\tableofcontents
\documentclass[11pt]{report}
\usepackage[utf8]{inputenc}
\usepackage[spanish]{babel}
\usepackage{amsmath,amssymb,amsthm,mathtools}
\usepackage{xcolor}
\usepackage{listings}
\usepackage{booktabs}
\usepackage{hyperref}
\usepackage[T1]{fontenc}
\usepackage{bussproofs}
\EnableBpAbbreviations
\let\AXM\AXC \let\BIM\BIC \let\UIM\UIC \let\TIM\TIC \let\RLM\RL

% --- Notación ---
\newcommand{\Imp}{\rightarrow}
\newcommand{\Not}{\neg}
\newcommand{\Or}{\vee}
\providecommand{\Sat}{\models}
\providecommand{\Prov}{\vdash}
\newcommand{\Provnd}{\vdash_{\mathrm{ND}}}
\newcommand{\sem}[1]{[\![ #1 ]\!]}

% --- Entornos de teoremas (declaración del autor) ---
\newtheorem{teorema}{Teorema}[chapter]
\newtheorem{lema}[teorema]{Lema}
\newtheorem{corolario}[teorema]{Corolario}
\newtheorem{proposicion}[teorema]{Proposición}
\theoremstyle{definition}
\newtheorem{definicion}[teorema]{Definición}
\newtheorem{ejemplo}[teorema]{Ejemplo}
\theoremstyle{remark}
\newtheorem{nota}[teorema]{Nota}
\newenvironment{demostracion}{\begin{proof}[Demostración]}{\renewcommand{\qedsymbol}{$\blacksquare$}\end{proof}}
\newenvironment{solucion}{\begin{proof}[Solución]}{\renewcommand{\qedsymbol}{$\square$}\end{proof}}

% --- Configuración de listings para Lean ---
\lstdefinelanguage{Lean}{
  morekeywords={import, namespace, end, inductive, def, theorem, abbrev,
    by, intro, apply, exact, simp, constructor, cases, with, match, open,
    have, let, refine, rcases, induction, noncomputable, attribute,
    instance, obtain, by_cases, by_contra, fun},
  sensitive=true, morecomment=[l]{--}, morestring=[b]"}
\lstset{
  language=Lean, basicstyle=\ttfamily\small,
  keywordstyle=\color{blue!60!black}, commentstyle=\color{green!40!black},
  stringstyle=\color{red!50!black}, showstringspaces=false,
  breaklines=true, frame=single, columns=fullflexible, mathescape=true,
  literate=
    {⟶}{{$\to$}}1 {~}{{$\neg$}}1 {∅}{{$\emptyset$}}1 {⊢}{{$\vdash$}}1
    {⊧}{{$\models$}}1 {⊨}{{$\models$}}1 {⋎}{{$\lor$}}1 {∼}{{$\sim$}}1
    {¬}{{$\neg$}}1 {∀}{{$\forall$}}1 {→}{{$\to$}}1 {↔}{{$\leftrightarrow$}}1
    {Γ}{{$\Gamma$}}1 {Δ}{{$\Delta$}}1 {φ}{{$\varphi$}}1 {ψ}{{$\psi$}}1
    {χ}{{$\chi$}}1 {∈}{{$\in$}}1 {∉}{{$\notin$}}1 {⊆}{{$\subseteq$}}1
    {∪}{{$\cup$}}1 {≠}{{$\neq$}}1 {⊥}{{$\bot$}}1 {∧}{{$\land$}}1
    {∨}{{$\lor$}}1 {·}{{$\cdot$}}1 {⟨}{{$\langle$}}1 {⟩}{{$\rangle$}}1
    {▸}{{$\triangleright$}}1 {á}{{\'a}}1 {…}{{$\dots$}}1}

\title{Completitud de la deducción natural proposicional en Lean~4}
\author{}
\date{}

\begin{document}
\maketitle
\tableofcontents
\input{ThesisReport_ND_ES.txt}
\end{document}

\end{document}

\end{document}

\end{document}
